\documentclass[12pt,a4paper]{article}
\usepackage[utf8]{inputenc}
\usepackage[margin=1in]{geometry}
\usepackage{graphicx}
\usepackage{listings}
\usepackage{xcolor}
\usepackage{booktabs}
\usepackage{array}
\usepackage{hyperref}

\lstset{
    basicstyle=\ttfamily\small,
    keywordstyle=\color{blue},
    commentstyle=\color{green!60!black},
    stringstyle=\color{red},
    breaklines=true,
    frame=single
}

\title{Smart Air Traffic Control Priority Landing Simulator}
\author{Fadi Muhammed\\B.Tech CSE (Cyber Security), Batch 2025-29\\Semester III}
\date{2025-2026}

\begin{document}

\maketitle

\section{Abstract}

This project presents a Smart Air Traffic Control (ATC) Priority Landing Simulator implemented in Java, integrating principles from Data Structures and Algorithms (DSA) and Object-Oriented Programming (OOP). The system manages landing requests from commercial, cargo, and emergency aircraft using custom data structures implemented from scratch without the Java Collections Framework. A priority-based queue determines landing order using a configurable scoring algorithm, while a circular queue models the holding pattern. Flight ID lookup uses a custom hash table with chaining, and binary search over a sorted registry enables efficient range queries. The design demonstrates inheritance, polymorphism, encapsulation, custom exception handling, and file persistence.

\textbf{Keywords:} Air Traffic Control, Priority Queue, Hash Table, Merge Sort, Binary Search, OOP, Java

\section{Introduction}

\subsection{Problem Statement}

Managing aircraft landing sequences at busy airports requires efficient prioritization considering multiple factors: fuel levels, emergency status, wait time, and aircraft type. Traditional ATC systems use manual sequencing, leading to inefficiencies. This project implements a menu-driven simulator for a small airport with limited runway capacity, where aircraft must be sequenced based on dynamic priority scores.

\subsection{Objectives}

\begin{itemize}
    \item Design class hierarchy using OOP concepts with proper encapsulation
    \item Implement at least four non-trivial data structures from scratch
    \item Apply sorting and searching algorithms with complexity analysis
    \item Demonstrate exception handling and file persistence
    \item Create user-friendly interface using JavaFX
\end{itemize}

\subsection{Scope}

The system models a single-runway airport with:
\begin{itemize}
    \item Flight registration, update, and cancellation
    \item Priority-based landing scheduling
    \item Holding pattern management
    \item Search by Flight ID and range queries
    \item Landing history logging with reports
\end{itemize}

\section{System Analysis and Design}

\subsection{Functional Requirements}

\begin{enumerate}
    \item Register, update, and cancel flight records
    \item Assign dynamic priority based on fuel, emergency status, wait time
    \item Manage holding pattern with circular queue (fixed capacity)
    \item Process landings in priority order
    \item Search flights by ID (hash table) and by time/priority range (binary search)
    \item Sort and display registry by ETA, fuel level, priority score
    \item Display holding queue and next-to-land aircraft
    \item Persist landing history to JSON file
    \item Generate reports: emergency landings, average wait time, runway utilization
    \item Exception handling for duplicate IDs, full holding pattern, invalid data
\end{enumerate}

\subsection{Non-Functional Requirements}

\begin{itemize}
    \item \textbf{Performance:} All operations must be $O(\log n)$ or better on average
    \item \textbf{Reliability:} Custom exceptions for meaningful error handling
    \item \textbf{Maintainability:} Modular design with package separation
    \item \textbf{Scalability:} Configurable capacities for all data structures
\end{itemize}

\section{Data Structures and Algorithms}

\subsection{1. Array-based Priority Queue}

Implemented using binary heap structure. Elements stored in array where parent-child relationships satisfy heap property based on priority score.

\textbf{Operations:}
\begin{itemize}
    \item Enqueue: $O(\log n)$ - bubble up
    \item Dequeue: $O(\log n)$ - bubble down
    \item Peek: $O(1)$
    \item Space: $O(n)$
\end{itemize}

\textbf{Justification:} Landing order determined by dynamic priority scores. Heap provides efficient retrieval of maximum priority element with minimal overhead.

\subsection{2. Circular Queue}

Array-based queue with circular indexing for holding pattern management.

\textbf{Operations:}
\begin{itemize}
    \item Enqueue/Dequeue: $O(1)$
    \item Peek: $O(1)$
    \item Space: $O(n)$
\end{itemize}

\textbf{Justification:} Fixed-capacity holding pattern requires FIFO behavior with constant-time operations. Circular array prevents unnecessary memory allocation.

\subsection{3. Hash Table with Chaining}

Custom hash table using division method for hash function and linked list chaining for collision resolution.

\textbf{Operations:}
\begin{itemize}
    \item Put/Get/Remove: $O(1)$ average, $O(n)$ worst case
    \item Space: $O(n)$
\end{itemize}

\textbf{Justification:} Flight ID lookup requires constant-time access. Hash table provides faster average access than binary search for single-key lookups.

\subsection{4. Sorted Array Registry}

Array maintained in sorted order by Flight ID, supporting binary search.

\textbf{Operations:}
\begin{itemize}
    \item Binary Search: $O(\log n)$
    \item Insert: $O(n)$ - requires shifting
    \item Range Query: $O(\log n + k)$ where $k$ is result size
    \item Space: $O(n)$
\end{itemize}

\textbf{Justification:} Range queries by Flight ID efficiently supported. Binary search significantly faster than linear search for large datasets.

\subsection{5. Stack for Undo Operations}

Linked-list based stack storing command history for cancellation/undo.

\textbf{Operations:}
\begin{itemize}
    \item Push/Pop/Peek: $O(1)$
    \item Space: $O(n)$
\end{itemize}

\subsection{6. Merge Sort}

Recursive divide-and-conquer sorting algorithm applied to registry views.

\textbf{Complexity:}
\begin{itemize}
    \item Time: $O(n \log n)$ best/average/worst
    \item Space: $O(n)$ auxiliary
\end{itemize}

\textbf{Justification:} Stable sort with guaranteed $O(n \log n)$ worst-case performance. Suitable for sorting by multiple criteria (ETA, fuel, priority).

\subsection{7. Binary Search}

Iterative binary search on sorted registry.

\textbf{Complexity:}
\begin{itemize}
    \item Time: $O(\log n)$
\end{itemize}

\textbf{Justification:} Efficient lookup for ordered data, forming foundation for range queries.

\section{Object-Oriented Design}

\subsection{Class Hierarchy}

\begin{lstlisting}
Aircraft (abstract)
├── CommercialFlight
├── CargoFlight
└── EmergencyFlight
\end{lstlisting}

\textbf{Aircraft} contains common attributes: flightId (final), origin, destination, ETA, fuelLevel, emergencyFlag, arrivalOrder. Abstract methods: \texttt{getPriorityScore()} and \texttt{getType()}.

\textbf{Inheritance:} Each subclass overrides priority calculation with type-specific weights.

\textbf{Polymorphism:} Dynamic dispatch ensures correct priority calculation per aircraft type.

\subsection{Package Structure}

\begin{lstlisting}
atc/
├── model/           # Aircraft hierarchy, Registry, History
├── datastructures/  # Custom DS implementations
├── exceptions/      # Custom checked exceptions
├── persistence/     # File I/O (JSON)
└── util/           # FlightType enum
\end{lstlisting}

\subsection{Exception Handling}

Custom checked exceptions:
\begin{itemize}
    \item \texttt{DuplicateFlightIDException}: Prevents duplicate registration
    \item \texttt{HoldingPatternFullException}: Manages capacity constraints
    \item \texttt{InvalidPriorityException}: Validates priority calculations
\end{itemize}

All exceptions provide meaningful messages with context data.

\section{Priority Scoring Algorithm}

Priority Score = $w_{emergency} \times emergencyFlag + w_{fuel} \times (100 - fuelLevel) + w_{wait} \times waitMinutes$

Default weights:
\begin{itemize}
    \item Commercial: Emergency=1000, Fuel=10, Wait=1
    \item Cargo: Emergency=800, Fuel=15, Wait=2
    \item Emergency: Emergency=10000, Fuel=0, Wait=0
\end{itemize}

Emergency flights always receive maximum priority, ensuring safety.

\section{Complexity Analysis}

\subsection{Data Structure Operations}

\begin{table}[h]
\centering
\begin{tabular}{l|c|c|c}
\textbf{Structure} & \textbf{Operation} & \textbf{Time} & \textbf{Space} \\
\hline
Priority Queue & Enqueue & $O(\log n)$ & $O(n)$ \\
 & Dequeue & $O(\log n)$ & \\
Circular Queue & Enqueue/Dequeue & $O(1)$ & $O(n)$ \\
Hash Table & Put/Get & $O(1)$ avg & $O(n)$ \\
Sorted Array & Binary Search & $O(\log n)$ & $O(n)$ \\
 & Insert & $O(n)$ & \\
Merge Sort & Sort & $O(n \log n)$ & $O(n)$ \\
\end{tabular}
\end{table}

\subsection{Overall System Complexity}

\begin{itemize}
    \item Register flight: $O(\log n)$ priority queue + $O(1)$ hash table + $O(n)$ sorted insert
    \item Process landing: $O(\log n)$ dequeue
    \item Search by ID: $O(1)$ hash lookup
    \item Range query: $O(\log n + k)$
\end{itemize}

\section{Implementation Details}

\subsection{Key Classes}

\textbf{FlightRegistry}: Coordinates all data structures, manages consistency between priority queue, hash table, and sorted registry.

\textbf{LandingHistory}: Maintains landing records with statistical analysis (average wait time, fuel at landing).

\textbf{FilePersistence}: Manual JSON parsing/writing without external libraries, demonstrating file I/O concepts.

\subsection{Sample Data}

Project includes \texttt{sample\_flights.json} with FlightRadar24-style data:
\begin{lstlisting}
{"flightId": "AI101", "type": "COMMERCIAL", 
 "origin": "BOM", "destination": "DEL", 
 "eta": "2025-01-15T14:30:00", "fuelLevel": 75, ...}
\end{lstlisting}

\section{Testing}

\subsection{Unit Tests}

\begin{itemize}
    \item \textbf{Model tests}: Aircraft creation, priority calculation, equals/hashCode
    \item \textbf{DS tests}: Stack push/pop, circular queue wraparound, heap property
    \item \textbf{Registry tests}: Registration, landing processing, holding pattern
    \item \textbf{File I/O tests}: JSON save/load, parse validation
\end{itemize}

All tests pass with expected outputs. Priority ordering verified: Emergency > Cargo (lower fuel) > Commercial.

\section{Results and Discussion}

\subsection{Performance Observations}

The integrated design satisfies all requirements:
\begin{itemize}
    \item Emergency flights processed immediately with 10000 priority score
    \item Holding pattern prevents runway overloading
    \item Hash table enables instant flight ID lookup
    \item Binary search supports efficient range queries
\end{itemize}

\subsection{Strengths}

\begin{itemize}
    \item All data structures implemented from scratch without Collections Framework
    \item Modular design enables easy extension
    \item Configurable priority weights allow customization
    \item Comprehensive exception handling
\end{itemize}

\subsection{Limitations}

\begin{itemize}
    \item Sorted array insertion is $O(n)$ due to shifting
    \item Memory usage higher than using built-in collections
    \item Manual JSON parsing lacks robustness of libraries
\end{itemize}

\section{Conclusion and Future Scope}

This project successfully demonstrates integration of DSA and OOP concepts in a real-world ATC simulation. The system efficiently manages aircraft landing priorities using custom data structures with appropriate complexity analysis.

\textbf{Future enhancements:}
\begin{itemize}
    \item JavaFX GUI with real-time queue visualization
    \item BST alternative to sorted array for $O(\log n)$ insertion
    \item Multi-runway simulation with graph-based conflict detection
    \item Weather-based dynamic priority adjustments
\end{itemize}

\section{References}

\begin{enumerate}
    \item Syllabus: 24CST201 Data Structures and Algorithms, 24CST203 OOP Java
    \item Question Paper: QP Code 24CCT201/2024/P/2
    \item FlightRadar24 API Documentation
\end{enumerate}

\end{document}
